\chapter[LITERATURE REVIEW]{Literature Review}
\label{ch:literature-review}

\section{Introduction}
\label{sec:lr-introduction}
% -----------------------------------------------------------------
% Short. Say how the review is organised and why. Pre-empt the
% obvious objection: much of the fusion evidence comes from remote
% sensing and autonomous driving rather than corrosion, and the
% domain mismatch is the motivation, not an oversight.
% One paragraph describing the search: databases, date window
% (2021 onward), original journal articles only.
% -----------------------------------------------------------------

\note{Literature review introduction to be drafted}


\section{Corrosion and Its Detection}
\label{sec:lr-corrosion}
% Draws on the Corrosion (15 rows) and Pitting Corrosion (8 rows)
% tabs. This is the section no computer-vision researcher could
% write, and the examiners are chemical engineers -- do not thin it.
% Mechanisms, morphology, why corrosion is visually confusable.

\note{Corrosion physics and detection to be drafted}


\section{Vision-Based Corrosion Detection}
\label{sec:lr-vision}
% RGB-only work: detection, classification and segmentation.
% From the Image Detection (14 col) and Image Segmentation tabs.
% End each subsection with what the modality cannot supply.

\note{RGB-based detection to be drafted}


\section{Complementary Sensing Modalities}
\label{sec:lr-modalities}

\subsection{Infrared Thermography}
% Active vs passive IRT; subsurface corrosion; the bimodal
% RGB + IR corrosion work that already exists and works.
\note{Thermography to be drafted}

\subsection{LiDAR and Geometric Sensing}
% LiDAR vs laser profilometry vs structured light -- the sensor-gap
% distinction that leg 2 of the problem statement rests on.
\note{LiDAR and geometric sensing to be drafted}

\subsection{Other Spectral Modalities}
% Hyperspectral. Include it here and state why it was not chosen,
% rather than leaving the omission for an examiner to find.
\note{Hyperspectral and alternatives to be drafted}


\section{Multimodal Sensor Fusion}
\label{sec:lr-fusion}
% From the Image Fusion (17 col) tab.
\subsection{Fusion Architectures}
% Early, intermediate and late fusion; what each costs and buys.
\note{Fusion architectures to be drafted}

\subsection{Registration and Calibration}
% Targetless extrinsic calibration of RGB + thermal + LiDAR is a
% solved problem in the literature -- this subsection is what makes
% the proposed system credible rather than speculative.
\note{Registration and calibration to be drafted}

\subsection{Evidence That the Added Modality Earns Its Place}
% The ablation evidence. This subsection is the spine of the whole
% review and feeds leg 3 directly. Report effect sizes honestly,
% including the modest ones.
\note{Ablation evidence to be drafted}


\section{Artificial Intelligence Methods}
\label{sec:lr-ai}
% From the Artificial Intelligence and CNN tabs. AI is the
% instrument, not a pillar of the research -- keep this section
% deliberately thin.
\note{AI methods to be drafted}


\section{Datasets and Benchmarks}
\label{sec:lr-datasets}
% Public corrosion datasets are RGB-only with one trimodal
% exception. State the data problem plainly here; the methodology
% chapter has to answer it.
\note{Datasets to be drafted}


\section{Research Gap}
\label{sec:lr-gap}
% Synthesis, not a summary. Three or four paragraphs that close the
% loop back to the Problem Statement in Chapter 1. Nothing new may
% appear here that was not evidenced above.
\note{Research gap synthesis to be drafted}
